%!%
% Copyright (c) 2026 mingcheng <mingcheng@outlook.com>
% 
% File: ch06-prompting-basics.tex
% Author: mingcheng <mingcheng@outlook.com>
% File Created: 2026-03-27 15:50:02
% 
% Modified By: mingcheng <mingcheng@outlook.com>
% Last Modified: 2026-03-29 07:50:18
%%

% !TeX root = ../prompt-engineering-guide-for-beginners.tex

\chapter{提示工程入门：最常用的四把钥匙}
\label{ch:prompting-basics}

前面几章我们打好了基础，现在正式进入提示工程的技巧部分。
这一章介绍的四种技巧是日常使用频率最高的，掌握它们就能覆盖你 80\% 以上的使用场景。

\section{零样本提示（Zero-Shot Prompting）}
\label{sec:zero-shot}

这是最简单直接的提示方式：不给任何示例，直接告诉 AI 你的需求。

\textit{类比：你跟一个聪明的实习生说「帮我翻译这段话」，他直接就能做，不需要你先示范。}

在机器学习领域，由于大语言模型在预训练阶段经过了海量任务数据的微调（Zero-Shot Transfer），它们已经获得了强大的“任务泛化（Generalization）”能力。所以对于很多人类常识范围内的任务，你不需要额外喂给它案例集，直接下达 System 级指令就好。

\textbf{示例：}

\begin{promptbox}
请把下面这段文字翻译成英文：

今天天气真好，适合出去走走。
\end{promptbox}

\begin{promptbox}
请将以下文字的情感分类为「正面」「负面」或「中性」：

这家餐厅的菜品味道一般，但服务态度还不错。
\end{promptbox}

零样本提示适用于\textbf{任务需求边界极度清晰、且无需在最终表现形式上做任何包装点缀的强功能性场景}。比如：单纯的中外文互译、常识科普型问答、文本的基础情绪归类、或者干脆把 AI 当成高级语病检查器来使用。

但很快你就会在实践中发现：当你期待 AI 的输出不仅要完成任务，还要带有特定的排版和风格（例如你公司特有的月报格式、或者某种特定的文风）时，零样本提示就不太够用了。这时候就需要用到下一种技巧：少样本提示。

\section{少样本提示（Few-Shot Prompting / In-Context Learning）}
\label{sec:few-shot}

有时候光靠语言描述很难讲清楚你到底想要什么样的输出。这时候最有效的办法是：直接给 AI 看几个例子，触发大模型的“上下文学习（In-Context Learning）”能力。

\textit{类比：你给实习生看了三份合格的周报样本，他立马知道格式、语气和重点该怎么写。}

\textbf{示例：}

\begin{promptbox}
请根据以下示例，为新的产品生成一句广告语：

产品：运动手环

广告语：24 小时守护你的每一次心跳

产品：降噪耳机

广告语：喧嚣世界里，给你一片宁静

产品：智能台灯

广告语：
\end{promptbox}

AI 会根据前面两个例子的模式，为「智能台灯」生成风格一致的广告语。

少样本提示的几个要点：

\begin{itemize}
  \item \textbf{样本数量}：通常提供 2 到 5 个示例即可。太少 AI 可能抓不住规律；太多则浪费上下文窗口的空间，甚至可能导致 AI 过度模仿而失去灵活性。
  \item \textbf{样本质量}：你提供的示例就是 AI 模仿的「标准答案」，所以务必确保每个示例本身在内容和表达上都是你满意的状态。
  \item \textbf{格式统一}：所有示例的结构、标点符号的使用方式、段落排版都应保持高度一致，像模板一样整齐。
  \item \textbf{覆盖多样性}：几个示例之间尽量涵盖不同的情况和类型，这样能帮助 AI 更好地理解规律并举一反三。
\end{itemize}

少样本提示特别适合以下场景：自定义格式的输出、特定风格的写作、分类任务需要定义类别的情况、以及你难以用语言精确描述期望输出的时候。

\section{链式思考提示（Chain-of-Thought, CoT）}
\label{sec:cot}

当你让 AI 处理需要推理的任务（比如数学题、逻辑判断、多步骤分析），直接问它要答案，它可能会给出错误的结果。但如果你让它「一步一步地想」，它的准确率会大幅提升。

\textit{类比：让实习生别直接报答案，先在纸上写出解题步骤。中间出错你也能看到并纠正。}

\textbf{不用链式思考的例子：}

\begin{promptbox}
小明有 15 个苹果，他给了小红 3 个，又从小华那里拿了 7 个。后来他把手里苹果的一半分给了妈妈。小明现在有几个苹果？
\end{promptbox}

AI 可能直接给一个答案，但不一定对。

\textbf{用了链式思考的例子：}

\begin{promptbox}
小明有 15 个苹果，他给了小红 3 个，又从小华那里拿了 7 个。后来他把手里苹果的一半分给了妈妈。小明现在有几个苹果？

请一步一步思考，展示你的推理过程。
\end{promptbox}

AI 会这样回答：
\begin{aibox}
第一步：小明开始有 15 个苹果。

第二步：给了小红 3 个，剩下 15 - 3 = 12 个。

第三步：从小华那里拿了 7 个，现在有 12 + 7 = 19 个。

第四步：把一半分给妈妈，19 / 2 = 9.5，取整为 9 个（如果不能分半个苹果的话）。

所以小明现在有 9 个半苹果（或者 9 个，取决于苹果能不能切开）。
\end{aibox}

使用链式思考的好处不仅是提高最终答案的准确率。因为模型展示了完整的推理过程，一旦某一步出现了逻辑错误，你就能迅速定位问题所在，然后在对话中要求它从出错的那一步重新推导。

使用链式思考提示的触发咒语极其简练，你甚至只需要在那些需要长距离逻辑折返的问题行末尾打上一句简单的：

\begin{itemize}
  \item 「请一步一步仔细思考」（泛用度极高的经典咏唱）
  \item 「请向我展示得出该推论的核心中间验证过程」
  \item 「请先对材料进行逐条剖析后再给出汇总结论」
  \item 「Let's think step by step」（玄学上的大招，这句英文在早先大基座模型实验效果上的提分异常显著，不妨一试）
\end{itemize}

链式思考提示适合所有需要逻辑推理的场景：数学计算、因果分析、利弊权衡、决策建议等。

\section{角色扮演提示（Role Playing）}
\label{sec:role-playing}

给 AI 指定一个角色身份，它的输出风格、专业度和关注点都会跟着变化。这个技巧简单但非常好用。

\textit{类比：告诉实习生「你现在假装是客户来投诉」，他的语气和关注点立刻不一样。}

\textbf{示例：}

\begin{promptbox}
你是一位拥有 20 年经验的资深营养师。我最近经常熬夜加班，饮食不规律，请给我一份适合上班族的一周健康饮食建议。
\end{promptbox}

\begin{promptbox}
你是一名专业的产品经理，擅长用户需求分析。请帮我分析以下用户反馈，提炼出三个最核心的需求：

（粘贴用户反馈内容）
\end{promptbox}

\begin{promptbox}
你是一位严厉但公正的面试官，请针对「Java 后端开发工程师」这个岗位，给我出 5 道技术面试题，并附上参考答案。
\end{promptbox}

角色扮演的一些技巧：

\begin{itemize}
  \item \textbf{角色越具体越好}：「资深营养师」比「营养方面的专家」更有效，「拥有 20 年经验的儿科医生」比「医生」更有效
  \item \textbf{可以叠加多个角色特征}：「你是一名擅长通俗化表达的科学家」，这样既有专业度又有可读性
  \item \textbf{可以让 AI 扮演你的对立面}：比如让它扮演「挑刺的客户」或「试图找漏洞的黑客」来帮你检查方案的薄弱环节
  \item \textbf{有趣的角色能激发创意}：让 AI 扮演「一个 5 岁小孩」来解释复杂概念，或扮演「毒舌评论家」来吐槽你的文案
\end{itemize}

这四种技巧（零样本、少样本、链式思考、角色扮演）是提示工程中使用频率最高的基础工具。在实际工作中，你完全可以把它们灵活组合使用。比如：

\begin{promptbox}
你现在是一位在华尔街工作多年的对冲基金经理（角色扮演），请一步一步分析以下几组备选方案的潜在风险（链式思考）。
\end{promptbox}

掌握了这四种基础技巧，你已经可以应对 90\% 以上的日常办公、文本处理和信息获取场景了。

\begin{tipbox}
快速选择指南：如果任务简单直接，先试零样本；如果 AI 理解不到位，加几个示例走少样本；如果需要严密推理，加上「请一步一步思考」开启链式思考；如果想要特定风格或专业视角，给 AI 分配一个角色。这四把钥匙可以单独使用，也可以自由组合。
\end{tipbox}

如果你还想让 AI 处理更复杂的任务，比如需要极高准确性的推理、多步骤协作等，下一章会介绍几种更进阶的增强技巧。
